\subsection{2. Những nghiên cứu liên quan}

\indent\indent Những năm gần đây, việc nhận diện tài khoản giả mạo (Sybil) trên mạng xã hội đang dịch chuyển mạnh mẽ từ các kỹ thuật khai thác đặc trưng phẳng sang phân tích cấu trúc mạng lưới. Al-Qurishi và cộng sự \cite{al2017sybil} đã phân loại các phương pháp tiếp cận thành hai nhóm chính: dựa trên đặc trưng hành vi, thông tin hồ sơ và dựa trên cấu trúc đồ thị xã hội. Nghiên cứu này đề cập hai dấu hiệu nhận biết các cụm Sybil về mặt hành vi: (1) \textit{sự đồng nhất về thời gian đăng ký} do kẻ tấn công sử dụng công cụ tự động để tạo lập hồ sơ hàng loạt; và (2) \textit{tính rập khuôn trong thông tin cá nhân}, nơi các thuộc tính như tên hiển thị hay tiểu sử thường được sinh ra từ các thuật toán chuỗi. Việc khai thác các dấu vết này cho phép nhận diện toàn bộ mạng lưới Sybil. Tuy nhiên, khi xét về mặt cấu trúc, các thuật toán đồ thị truyền thống như SybilRank lại thường bỏ qua các đặc trưng dữ liệu này và chỉ dựa trên giả định rằng tài khoản giả mạo khó có thể thiết lập số lượng lớn kết nối với người dùng thật. Giả định này hiện không còn chính xác trong kỷ nguyên Web3, nơi các kẻ tấn công sử dụng mã lệnh tự động để tạo ra hàng nghìn liên kết chéo nhằm trục lợi từ các chương trình tặng thưởng (Airdrop) \cite{liu2025detecting}.\vspace{0.3cm}

Để giải quyết bài toán biểu diễn dữ liệu phức tạp, mạng nơ ron đồ thị (GNN) đã trở thành hướng đi chủ đạo. Zhang và cộng sự \cite{zhang2019graph} đã chứng minh hiệu quả của mạng GCN trong việc học biểu diễn nút, nhưng đồng thời cũng chỉ ra hiện tượng làm mịn quá mức trên đồ thị siêu đặc. Khắc phục điểm yếu này, Veličković và cộng sự \cite{velickovic2017graph} đề xuất mạng GAT, cho phép mô hình gán trọng số khác nhau cho từng láng giềng. Nhằm tối ưu hóa hơn nữa, Brody và cộng sự \cite{brody2021how} đã giới thiệu phiên bản cải tiến GATv2, khắc phục hạn chế của cơ chế chú ý tĩnh ở bản gốc để mang lại khả năng chú ý động với sức mạnh biểu diễn vượt trội. Mới đây, Heeb và cộng sự \cite{heeb2024sybil} đã ứng dụng GAT vào phát hiện Sybil, khẳng định ưu thế của cơ chế chú ý trong việc định tuyến tín hiệu dị thường. Song song đó, việc trích xuất đặc trưng văn bản từ tiểu sử người dùng thông qua các mô hình ngôn ngữ như Sentence-BERT \cite{reimers2019sentence} cũng được kết hợp để biểu diễn đặc trưng cho các nút trên đồ thị.\vspace{0.3cm}

Mặc dù GNN mang lại hiệu năng cao, thách thức vẫn là sự thiếu hụt dữ liệu nhãn thực tế. Để lấp đầy khoảng trống này, các phương pháp học không giám sát như VGAE của Kipf và Welling \cite{kipf2016variational} mở ra tiềm năng trích xuất đặc trưng nhúng mà không cần nhãn. Đồng thời, nhu cầu về tính giải thích (Explainability) cũng trở nên cấp thiết. Mặc dù các công cụ như GNNExplainer \cite{ying2019gnnexplainer} đã xuất hiện, nhưng chi phí tính toán lớn khiến chúng khó ứng dụng rộng rãi. Hơn nữa, việc khai thác các đặc trưng cạnh như quyền đồng sở hữu ví vẫn chưa được các mô hình GNN tiêu chuẩn xử lý tối ưu, dù Gong và Cheng \cite{gong2019exploiting} đã nhấn mạnh tầm quan trọng của chúng.\vspace{0.3cm}

Tóm lại, khoảng trống nghiên cứu hiện tại nằm ở ba điểm: (1) \textit{Thiếu một quy trình khép kín giải quyết bài toán thiếu nhãn trong Web3;} (2) \textit{Chưa tận dụng tối đa các ràng buộc về tính sở hữu đặc thù của SocialFi;} và (3) \textit{Sự đánh đổi giữa hiệu năng của mô hình học sâu và tính minh bạch trong diễn giải kết quả rủi ro.} Đây chính là động lực để hệ thống \textbf{SybilSignal} ra đời, sử dụng kiến trúc đồ thị đa tầng nhận thức ràng buộc và kiến trúc phân lớp tách rời (Decoupled ML) để tối ưu hóa đồng thời cả độ chính xác và khả năng giải thích.